\chapter{Análisis de Pequeña Señal y Parámetros del Amplificador}
\label{chap:pequena_senal}

En este capítulo se desarrolla el modelo linealizado de pequeña señal para la configuración en emisor común a partir
de los parámetros híbridos del transistor bipolar, deduciendo analíticamente las expresiones de ganancia y de
impedancia. Asimismo, se describen los fundamentos de los métodos de medición experimental utilizados en laboratorio.

\section{Modelo Híbrido Equivalente del Transistor}

Para señales de amplitud reducida superpuestas al punto de reposo $Q$, el transistor BJT se comporta como un
dispositivo lineal. Despreciando los efectos de realimentación inversa ($h_{re} \approx 0$) y la conductancia de
salida ($h_{oe} \approx 0$) frente a la carga de colector, el transistor queda representado por su resistencia
dinámica de entrada $h_{ie}$ y su generador de corriente controlado por corriente $h_{fe} \cdot i_b$.

En la figura~\ref{crkt:pequena_senal} se presenta el circuito equivalente de pequeña señal del amplificador completo.

\begin{figure}[!ht]
  \centering
  \begin{circuitikz}[scale=0.9, transform shape]
  % Terminal de entrada
  \draw (0,0) node[ground] {}
    to[open, v^>=$v_i$] (0,2.5)
    to[short, o-*] (1.5,2.5) coordinate (in_node);

  % Resistencia RB = R1 // R2
  \draw (in_node) to[R, l=$R_B$, -] (1.5,0) node[ground] {};

  % Modelo híbrido del transistor
  \draw (in_node) to[short] (3.0,2.5)
    to[R, l=$h_{ie}$, i=$i_b$] (3.0,0) node[ground] {};

  % Fuente dependiente de corriente de colector
  \draw (5.5,0) node[ground] {}
    to[cisourceAM, l=$h_{fe} i_b$] (5.5,2.5)
    to[short, -] (7.0,2.5) coordinate (c_node);

  % Resistencia de colector RC
  \draw (7.0,2.5) to[R, l=$R_C$, -] (7.0,0) node[ground] {};

  % Resistencia de carga RL
  \draw (c_node) to[short] (8.5,2.5) coordinate (l_node)
    to[R, l=$R_L$, -] (8.5,0) node[ground] {};

   % 1. Dibujamos el terminal de salida
  \draw (l_node) to[short, -o] (9.5,2.5) node[above] {$v_L$} coordinate (out_node);

  % 2. Colocamos la masa en el punto inferior y medimos la tensión en abierto hacia arriba
  \draw (9.5,0) node[ground] {} to[open, v_>=$v_L$] (out_node);
\end{circuitikz}

  \caption{Modelo equivalente de pequeña señal en parámetros híbridos para el amplificador emisor común.}
  \label{crkt:pequena_senal}
\end{figure}

A partir de la física del semiconductor, la resistencia dinámica de la juntura base-emisor polarizada en directa es:
\begin{equation}
  r_e = \frac{V_T}{I_{CQ}}
  \label{eq:re_prima}
\end{equation}
donde $V_T = k T / q \approx 25.0\mV$ representa el potencial térmico a temperatura ambiente de laboratorio ($T \approx 290\Kelvin$).

La resistencia dinámica de entrada del transistor en emisor común ($h_{ie}$) resulta:
\begin{equation}
  h_{ie}  \approx h_{fe} \cdot \frac{V_T}{I_{CQ}} = 284 \cdot \frac{25.0\mV}{7.683\mA} \approx 924.08\ohm
  \label{eq:hie_calculo}
\end{equation}
y la ganancia de corriente en cortocircuito es $h_{fe} = \beta = 284$.

\section{Cálculo Analítico de Parámetros del Cuadripolo}

A partir del análisis de mallas y nodos del circuito de la figura~\ref{crkt:pequena_senal}, se deducen las expresiones
de los parámetros fundamentales del amplificador.

\subsection{Impedancia de Entrada (\texorpdfstring{$Z_i$}{Zi})}
La impedancia que la etapa presenta a la fuente de excitación comprende el paralelo de la red de polarización de base
($R_B = R_1 \parallel R_2$) y la resistencia de entrada del BJT ($h_{ie}$):
\begin{equation}
  Z_i = R_1 \parallel R_2 \parallel h_{ie} = R_B \parallel h_{ie}
  \label{eq:Zi_formula}
\end{equation}

Evaluando con $R_B = 5790.39\ohm$ y $h_{ie} = 924.08\ohm$:
\begin{equation}
  Z_i = \frac{5790.39\ohm \cdot 924.08\ohm}{5790.39\ohm + 924.08\ohm} \approx 796.90\ohm
  \label{eq:Zi_teorica}
\end{equation}

\subsection{Impedancia de Salida (\texorpdfstring{$Z_o$}{Zo})}
La impedancia de salida se evalúa anulando la fuente de excitación de entrada ($v_i = 0$). En dicha condición, la
corriente de base se extingue ($i_b = 0$), lo cual anula el generador dependiente de colector ($h_{fe} i_b = 0$).
Por lo tanto, la impedancia de salida vista desde los terminales de la carga resulta:
\begin{equation}
  Z_o = R_C = 1200\ohm
  \label{eq:Zo_teorica}
\end{equation}

\subsection{Ganancia de Tensión (\texorpdfstring{$A_v$}{Av})}
La tensión alterna desarrollada sobre la carga $R_L$ viene dada por la corriente inyectada por el colector sobre el
paralelo de resistencias:
\begin{equation}
  v_L = -h_{fe} \cdot i_b \cdot (R_C \parallel R_L)
\end{equation}
mientras que la tensión en los terminales de entrada es:
\begin{equation}
  v_i = i_b \cdot h_{ie}
\end{equation}

El cociente entre ambas tensiones define la ganancia de tensión de la etapa:
\begin{equation}
  A_v = \frac{v_L}{v_i} = -\frac{h_{fe} \cdot (R_C \parallel R_L)}{h_{ie}} = -\frac{R_C \parallel R_L}{r_e'}
  \label{eq:Av_formula}
\end{equation}

Reemplazando los valores numéricos:
\begin{equation}
  A_v = -\frac{545.45\ohm}{3.254\ohm} \approx -167.62
  \label{eq:Av_teorica}
\end{equation}
donde el signo negativo explicita la inversión de fase de $180^\circ$ de la topología emisor común.

\subsection{Ganancia de Corriente (\texorpdfstring{$A_i$}{Ai})}
La ganancia de corriente se define como la razón entre la corriente alterna que atraviesa la carga ($i_L$) y la
corriente total que suministra el generador a la entrada del amplificador ($i_i$):
\begin{equation}
  i_L = \frac{v_L}{R_L}, \quad i_i = \frac{v_i}{Z_i}
\end{equation}
\begin{equation}
  A_i = \frac{i_L}{i_i} = \frac{v_L / R_L}{v_i / Z_i} = A_v \cdot \frac{Z_i}{R_L}
  \label{eq:Ai_formula}
\end{equation}

En valor absoluto:
\begin{equation}
  |A_i| = 167.62 \cdot \frac{796.90\ohm}{1000\ohm} \approx 133.58
  \label{eq:Ai_teorica}
\end{equation}

\section{Métodos de Medición Experimental en Laboratorio}

La determinación práctica de las impedancias en pequeña señal no puede realizarse de forma directa con un ohmímetro
convencional, ya que los instrumentos de corriente continua alterarían los puntos de polarización del transistor. Por
ello, se recurre al método de inserción de una resistencia serie ($R_S$) calibrada.

\subsection{Medición de \texorpdfstring{$Z_i$, $A_v$ y $A_i$}{Zi, Av y Ai} (Método del Resistor Serie de Entrada)}
En la figura~\ref{crkt:medicion_zi} se ilustra el esquema de conexionado.

\begin{figure}[!ht]
  \centering
  \begin{circuitikz}[scale=0.9, transform shape]
  % Generador y resistencia serie
  \draw (0,0) node[ground] {}
    to[sV, l=$v_s$] (0,2.5)
    to[R, l=$R_S$, *-] (2.2,2.5) coordinate (vi_node)
    to[short, -o] (2.2,2.5);

  \draw (0,0) to[short] (2.2,0) ;
  \draw (0,2.5) to[short, *-o] (0,3.2) node[above] {CH1 ($v_s$)};
  \draw (vi_node) to[short, *-o] (2.2,3.2) node[above] {CH2 ($v_i$)};

  % Bloque amplificador
  \draw (2.7, -0.5) rectangle (5.5, 3.0);
  \node at (4.1, 1.25) [align=center] {Amplificador\\Emisor Común};

  \draw (vi_node) to[short] (2.7, 2.5);
  \draw (2.2, 0) to[short] (2.7, 0);

  % Carga y medición de salida
  \draw (5.5, 2.5) to[short, -*] (6.8, 2.5) coordinate (vl_node)
    to[R, l=$R_L$] (6.8, 0) node[ground] {};
  \draw (5.5, 0) to[short] (6.8, 0);
  \draw (vl_node) to[short, -o] (7.5, 2.5) node[right] {$v_L$ (Salida)};
\end{circuitikz}

  \caption{Esquema de medición de $Z_i$, $A_v$ y $A_i$ intercalando un resistor serie $R_S$ en la entrada.}
  \label{crkt:medicion_zi}
\end{figure}

Se intercala un resistor conocido $R_S = 820\ohm$ entre el generador de funciones y la entrada del amplificador.
Mediante un osciloscopio de dos canales se miden simultáneamente la tensión suministrada por el generador ($v_s$) y la
tensión efectiva que ingresa a la base tras la caída en $R_S$ ($v_i$).

La corriente que ingresa a la etapa resulta de la ley de Ohm sobre el resistor testigo:
\begin{equation}
  i_i = \frac{v_s - v_i}{R_S}
  \label{eq:ii_exp}
\end{equation}

Conociendo $i_i$ y $v_i$, la impedancia de entrada se determina de forma inmediata como:
\begin{equation}
  Z_i = \frac{v_i}{i_i} = \frac{v_i \cdot R_S}{v_s - v_i}
  \label{eq:Zi_exp_formula}
\end{equation}

Simultáneamente, midiendo la amplitud pico a pico de la tensión sobre la carga $R_L$ ($v_L$), se computan la ganancia
de tensión experimental:
\begin{equation}
  A_v = \frac{v_L}{v_i}
  \label{eq:Av_exp_formula}
\end{equation}
y la ganancia de corriente experimental:
\begin{equation}
  A_i = \frac{i_L}{i_i} = \frac{v_L / R_L}{(v_s - v_i) / R_S}
  \label{eq:Ai_exp_formula}
\end{equation}

\subsection{Medición de \texorpdfstring{$Z_o$}{Zo} (Método de Inyección en Salida)}
Para evaluar la impedancia de salida de la etapa, se utiliza el montaje de la figura~\ref{crkt:medicion_zo}.

\begin{figure}[!ht]
  \centering
  \begin{circuitikz}[scale=0.9, transform shape]

  % Bloque amplificador
  \draw (1.2, -0.5) rectangle (4.0, 3.0);
  \node at (2.6, 1.25) [align=center] {Amplificador\\Emisor Común};

  \draw (0.0, 2.0) to[short] (1.2, 2.0);
  \draw (0.0, 0) to[short] (1.2, 0);
  \draw (0.0, 0) to[short] (0, 2);
  \draw (0.6, 0) node[ground] {} ;

  % Red de inyección en salida
  \draw (4.0, 2.0) to[short, -*] (5.2, 2.0) coordinate (vo_node)
    to[R, l_=$R_S$] (7.2, 2.0)
    to[sV, l_=$v_s$] (7.2, 0) node[ground] {};
  \draw (4.0, 0) to[short] (7.2, 0);

  \draw (vo_node) to[short, *-o] (5.2, 3.0) node[above] {CH2 ($v_o$)};
  \draw (7.2, 2.0) to[short, *-o] (7.2, 3.0) node[above] {CH1 ($v_s$)};
\end{circuitikz}

  \caption{Esquema de inyección de señal para la medición de la impedancia de salida $Z_o$.}
  \label{crkt:medicion_zo}
\end{figure}

Se cortocircuita a masa la entrada de señal en alterna ($v_{in} = 0$) para anular cualquier excitación propia del
transistor, conservando intacta su polarización en continua. En el terminal de colector se desacopla la carga $R_L$ y se
conecta el generador de funciones en serie con un resistor patrón de $R_S = 1200\ohm$.

Midiendo la tensión del generador ($v_s$) y la tensión en el colector ($v_o$), la corriente inyectada es:
\begin{equation}
  i_o = \frac{v_s - v_o}{R_S}
\end{equation}
y la impedancia de salida queda determinada como:
\begin{equation}
  Z_o = \frac{v_o}{i_o} = \frac{v_o \cdot R_S}{v_s - v_o}
  \label{eq:Zo_exp_formula}
\end{equation}
